\section{Theoretical Framework for Linear Decoding}
\label{sec:setup-linear-decoding}



We use $\mbX \in \reals^{M \times N_X}$ and $\mbY \in \reals^{M \times N_Y}$ to denote matrices holding responses to $M$ stimulus conditions measured across neural populations consisting of $N_X$ and $N_Y$ neurons, respectively.
Throughout this paper we will assume that the neural responses have been preprocessed so that the columns of $\mbX$ and $\mbY$ have zero mean.

\subsection{Linear decoding from neural population responses}

We first consider the problem of predicting (or ``decoding'') a target vector $\mbz \in \reals^M$ from neural population responses by a linear function $\mbX \mapsto \mbX \mbw$ where $\mbw \in \reals^{N_X}$.
One can view this as a simplified neural circuit model where each element of $\hat{\mbz} = \mbX \mbw + \mb{\epsilon}$ is the firing rate readout unit the $M$ conditions.
Here, we interpret $\mbw$ as a vector of synaptic weights and the random vector $\mb{\epsilon}$ represents potential noise corruptions; for example, $\mb{\epsilon}_i \sim \cN(0, \sigma^2)$ where $\sigma^2 > 0$ specifies the scale of noise.
This neural circuit interpretation of linear decoding is fairly standard within the literature~\cite{fusi2016neurons,kaufman2014cortical}.

Within this setting, we formalize the problem of decoding $\mbz$ from the population activity $\mbX$ through the following class of optimization problems:
\begin{align}
\label{eq:readout-optimization-class}
\underset{\mbw}{\text{maximize}} \quad \tfrac{1}{M} \mbz^\top \mbX \mbw - \tfrac{1}{2} \mbw^\top \mbG(\mbX) \mbw
\end{align}
where $\mbG(\cdot)$ is a function mapping $\reals^{M \times N}$ onto symmetric positive definite $N \times N$ matrices.
The objective is concave, and equating the gradient to zero yields a closed form solution:
\begin{equation}
\label{eq:linear-readout-closed-form}
\mbw^* = \tfrac{1}{M} \mbG(\mbX)^{-1} \mbX^\top \mbz
\end{equation}
To appreciate the relevance of \cref{eq:readout-optimization-class}, note that if $\mbG(\mbX) =  \mbC_X + \lambda \mbI$, where $\mbC_X:=\tfrac1M\mbX^\top\mbX$ is the empirical covariance of $\mbX$, then the solution $\mbw^*$ coincides with the optimum of the familiar ridge regression problem:
\begin{equation}
\label{eq:ridge}
\underset{\mbw}{\text{minimize}} \quad \tfrac{1}{M} \Vert \mbz - \mbX \mbw \Vert_2^2 + \lambda \Vert \mbw \Vert_2^2
\end{equation}
with solution $\mbw^* = \frac{1}{M}(\mbC_X + \lambda \mbI)^{-1} \mbX^\top\mbz$, where $\lambda > 0$ specifies the strength of regularization on the coefficients.
More generally, we can interpret \cref{eq:readout-optimization-class} as maximizing the inner product between the target, $\mbz$, and the linear readout, $\mbX \mbw$, plus a penalty term on the scale of $\mbw$ as quantified by the function $\mbw \mapsto \sqrt{\mbw^\top \mbG(\mbX) \mbw}$, which is a norm on $\reals^N$.

\Cref{eq:readout-optimization-class} is a useful generalization of \cref{eq:ridge} because there are situations where the scale of the decoded signal matters.
For instance, when the readout unit is corrupted with noise, $\hat{\mbz} = \mbX \mbw + \mb{\epsilon}$ where $\mb{\epsilon}_i \sim \cN(0, \sigma^2)$, we can interpret \cref{eq:readout-optimization-class} as maximizing the signal-to-noise ratio subject to a cost on the magnitude of the  weights $\mbw$.
% Moreover, if we use $\mbz_i \in \{-1, +1\}$ to encode binary class labels, the \cref{eq:readout-optimization-class} is essentially equivalent to fitting a support vector machine (SVM) to the neural responses \hl{cite}.
% That is, the solution (assuming one exists) to:
% \begin{equation}
% \label{eq:svm-classifier}
% \begin{aligned}
% & \underset{x}{\text{minimize}}
% & & \mbw^\top \mbG(\mbX) \mbw \\
% & \text{subject to}
% & & \mbz^\top \mbX \mbw \geq c
% \end{aligned}
% \end{equation}
% will coincide with \cref{eq:readout-optimization-class} for some choice of $c > 0 $.
% This can be seen by re-writing the problem in terms of Langrange multipliers and ignoring constant terms.


In most cases, we will consider choices of $\mbG(\cdot)$ that can be written as:
\begin{equation}
\label{eq:G-alpha-beta}
\mbG(\mbX) = a \mbC_X + b \mbI
\end{equation} 
for some $a,b \geq 0$.
Under this choice, the penalty on $\mbw$ equals the sum of two terms:
\begin{equation}\label{eq:G-a-b}
\mbw^\top \mbG(\mbX) \mbw = \frac{a}{M} \Vert \mbX \mbw \Vert_2^2 + b \Vert \mbw \Vert_2^2.
\end{equation} 
% which can each be interpreted as constraints on the neural circuit model diagrammed in \cref{fig:readouts_cartoon}.
The first term, scaled by $a$, penalizes the activation level of the readout unit, which can be viewed as an energetic constraint.
The second term, scaled by $b$, penalizes the strength of the synaptic weights, which can be viewed as a resource constraint for building and maintaining strong synapses.
We note again that ridge regression is achieved as a special case when $a = 1$ and $b = \lambda$.

\subsection{Quantifying differences between networks through the lens of decoding}



\begin{figure}
    \centering
    \includegraphics[scale=0.25] %[width=12.9cm]  %0.9\textwidth]
    {readouts_simple.pdf}\caption{Schematic of the proposed framework for comparing representations $\mbX$ and $\mbY$ (each dot represents mean neural responses to one of $M$ conditions) in terms of a decoding target $\mbz$. First, optimal linear decoding weights $\mbw^*$ and $\mbv^*$ are computed. Then the similarity between the two systems is measured in terms of the \textit{decoding similarity}: $\langle \mbX \mbw^*, \mbY \mbv^* \rangle$.}
    \label{fig:readouts_simple} 
\end{figure}

We now turn our attention to the the problem of quantifying representational similarity between $\mbX \in \reals^{M \times N_X}$ and $\mbY \in \reals^{M \times N_Y}$.
As before, we specify a target vector $\mbz \in \reals^M$ and optimize $\mbw$ as specified in \cref{eq:readout-optimization-class}.
We optimize $\mbv$ accordingly:
\begin{align}
\underset{\mbv}{\text{maximize}} \quad \tfrac{1}{M} \mbz^\top \mbY \mbv - \tfrac{1}{2} \mbv^\top \mbG(\mbY) \mbv
\end{align}
yielding the solution $\mbv^* = \tfrac{1}{M} \mbG(\mbY)^{-1} \mbY^\top \mbz$.
We are interested in quantifying similarity between $\mbX$ and $\mbY$ by comparing the behavior of these optimal decoders.
A straightforward way to quantify the alignment between the decoded signals is to compute their inner product, which we call the \emph{decoding similarity} (\cref{fig:readouts_simple}):
\begin{equation}
\label{eq:inner-prod-linear-optimal-readouts-fixed-z}
\langle \mbX \mbw^*, \mbY \mbv^* \rangle = \mbw^{*\top} \mbX^\top \mbY \mbv^* = \mbz^\top \mbK_X\mbK_Y \mbz 
\end{equation}
where we have leveraged \cref{eq:linear-readout-closed-form} and the analogous closed form expression for $\mbv^*$ while introducing the following normalized kernel similarity matrices:
\begin{equation}
\label{eq:change-of-variables-generalized-whitening}
    \mbK_X:=\frac{1}{M}\mbX\mbG(\mbX)^{-1}\mbX^\top \quad \text{and} \quad \mbK_Y:=\frac{1}{M}\mbY\mbG(\mbY)^{-1}\mbY^\top.
\end{equation}
Note that we have suppressed the dependence of $\mbK_X$ and $\mbK_Y$ on the function $\mbG(\cdot)$ to reduce notational clutter.

\Cref{eq:inner-prod-linear-optimal-readouts-fixed-z} is a reasonable quantification of similarity between $\mbX$ and $\mbY$ with respect to a fixed decoding task specified by $\mbz \in \reals^M$, and it may be a fruitful approach for hypothesis-driven comparisons of neural systems. We comment further on this possibility in our discussion.
On the other hand, many neural representations support a variety of downstream behavioral tasks.
For example, features extracted in early stages of visual processing (edges and textures) can be used to support many visual tasks such as object classification, segmentation, image compression, et cetera.
How can \cref{eq:inner-prod-linear-optimal-readouts-fixed-z} be adapted to quantify the similarity between $\mbX$ and $\mbY$ across more than one pre-specified decoding task?
We explore three simple ideas.

\textbf{Option 1 --- best case alignment.}
An optimistic measure of representational similarity between networks is to find the decoding task $\mbz \in \reals^M$ that results in a maximal decoder alignment.
Formally,
\begin{equation}
\label{eq:best-case-alignment}
\max_{\Vert \mbz \Vert_2 = 1} \langle \mbX \mbw^*, \mbY \mbv^* \rangle = \max_{\Vert \mbz \Vert_2 = 1} \mbz^\top \mbK_X \mbK_Y \mbz
\end{equation}
where the constraint that $\Vert \mbz \Vert_2 = 1$ is needed to keep the maximal value finite.

\textbf{Option 2 --- worst case alignment.}
An alternative is to find the  task that minimizes alignment:
\begin{equation}
\label{eq:worst-case-alignment}
\min_{\Vert \mbz \Vert_2 = 1} \langle \mbX \mbw^*, \mbY \mbv^* \rangle = \min_{\Vert \mbz \Vert_2 = 1} \mbz^\top \mbK_X \mbK_Y \mbz
\end{equation}
which is obviously the pessimistic counterpart to option 1 above.

\textbf{Option 3 --- average case alignment.}
Instead of maximizing or minimizing over $\mbz$, we can approach the problem by averaging over a distribution of decoding tasks.
Formally, we can quantify alignment in terms of the \textit{average decoding similarity}:
\begin{equation}
\label{eq:average-case-alignment}
\expect \langle \mbX \mbw^*, \mbY \mbv^* \rangle
\end{equation}
where the expectation is taken with respect to $\mbz \sim P_{\mbz}$.

In the next subsection, we show that the third option is closely related to several existing neural representational similarity measures.
However, all three quantities are potentially of interest and they are easy to compute as we document below in \cref{proposition:best-worst-case,proposition:avg-case}.
% Deriving these results is straightforward, so we provide proofs below.

\begin{proposition}
\label{proposition:best-worst-case}
The best case and worst case decoding alignment scores, \cref{eq:best-case-alignment,eq:worst-case-alignment}, are respectively given by the largest and smallest eigenvalues of:
\begin{equation}
\frac{1}{2} \left ( \mbK_X \mbK_Y + \mbK_Y \mbK_X \right )
\end{equation}
\end{proposition}

\begin{proposition}
\label{proposition:avg-case}
The average decoding similarity, \cref{eq:average-case-alignment}, is given by:
\begin{align}
\label{eq:proposition-avg-case}
\expect \langle \mbX \mbw^*, \mbY \mbv^* \rangle =
\Tr \left( \mbK_X \mbK_z \mbK_Y \right),
\end{align}
where $\mbK_z:=\E[\mbz\mbz^\top]$ represents a kernel matrix for the targets.
\end{proposition}

\begin{proof}[Proof of \cref{proposition:best-worst-case}]
Let $\mbA \in \reals^{M \times M}$ be any matrix, potentially non-symmetric.
We can express $\mbA$ as a sum of a symmetric and skew symmetric matrix, $\mbA = \mbB + \mbC$, where:
\begin{equation}
\mbB = \frac{1}{2}\left ( \mbA + \mbA^\top \right) \qquad \text{and} \qquad \mbC = \frac{1}{2}\left ( \mbA - \mbA^\top \right)
\end{equation}
For any $\mbz \in \reals^M$, we have $\mbz^\top \mbA \mbz = \mbz^\top \mbB \mbz + \mbz^\top \mbC \mbz$.
It is easy to verify that $\mbz^\top \mbC \mbz$ equals zero.
Thus, $\mbz^\top \mbA \mbz = \mbz^\top \mbB \mbz$, and maximizing over $\mbz$ subject to $\Vert \mbz \Vert_2$ yields the largest eigenvalue of $\mbB$ (since it is symmetric).
Likewise, minimizing yields the smallest eigenvalue of $\mbB$.
Substituting $\mbA = \mbK_X \mbK_Y$ into this argument proves the proposition.
\end{proof}

\begin{proof}[Proof of \cref{proposition:avg-case}]
Using \cref{eq:inner-prod-linear-optimal-readouts-fixed-z}, we have:
\begin{equation*} 
\expect \langle \mbX \mbw^*, \mbY \mbv^* \rangle = \expect [\mbz^\top \mbK_X \mbK_Y \mbz ].
\end{equation*}
Using the cyclic property of the trace operator to manipulate this further, we get:
\begin{equation*}
\mbz^\top \mbK_X \mbK_Y \mbz  = \Tr( \mbK_X\mbz\mbz^\top\mb \mbK_Y).
\end{equation*}
Taking the expectation with respect to $\mbz$ and using linearity of trace yields \cref{eq:proposition-avg-case}.
\end{proof}




